\input epsf
\input eplain
\input miniltx %
\input graphicx.sty %


% Example of setting red text in plain TeX using pdfTeX
%\pdfoutput=1 % Ensure PDF output
%\def\red{\pdfcolorstack0 push {1 0 0 rg 1 0 0 RG}}
%\def\black{\pdfcolorstack0 pop}
%\def\green{\pdfcolorstack0 push {0 1 0 rg 0 1 0 RG}}

% Bold and italic bold scaled fonts for sections and examples.
\font\small         = cmr7
\font\sxmplbx       = cmbx10 
\font\sxmplbxti     = cmbxti10 
\font\xmplbx        = cmbx12 % scaled \magstephalf
\font\mxmplbx       = cmbx12 at 16pt
\font\bxmplbx       = cmti12 at 18pt %logo1%0cmbx12 scaled \magstep2
\font\bxmplbxt      = cmti12 at 20pt %logo1%0cmbx12 scaled \magstep2
\font\chap          = cmti14
\font\chapcont      = cmti12
\font\sectioncont   = cmti10
\font\header        = cmti10

%\font\tenrm         = pplr at 12pt
%\font\tentt         = ppls at 12pt
%\font\tenbf         = pplb at 12pt
%\font\tenit         = pplri at 12pt
%\font\header        = ptmrc at 12pt

%\font\tenrm         = pxr at 12pt
%\font\tentt         = cmtt10 at 12pt
%\font\teni          = pxmi at 12pt
%\font\tenx          = pxex at 12pt
%\font\tenbf         = pxb at 12pt
%\font\tenit         = pxi at 12pt
%\font\tensl         = pxsl at 12pt
%\font\header        = ptmrc at 12pt

%\font\smallmath = pplri at 12pt
%\font\xmplbxti  = cmbxti10 scaled \magstephalf
\font\xmplbxti  = cmbxti10 at 14pt
\font\bookfont  = cmr11

%\font\smallmath = cmu10
%\font\xmplbxti  = cmbxti10 scaled \magstephalf

\newdimen\leftfigureskip \leftfigureskip=10pc
\newdimen\rightfigureskip \rightfigureskip=5pc

% Set the font to Palentino
% NOTE: the font switches must keep their \fam assignments (as in plain
% TeX's own definitions).  A bare \let\rm=\tenrm changes only the current
% text font; in math mode letters render from the math family, so
% \sin, \cos, \log, ... (which expand to \mathop{\rm ...}) would silently
% come out in math italic.
\def\pal{\def\rm{\fam0 \tenrm}%
   \def\it{\fam\itfam \tenit}%
   \def\bf{\fam\bffam \tenbf}%
   \def\tt{\fam\ttfam \tentt}%
   \let\mit=\teni
   \let\mex=\tenx
   \rm
   }%

% Table of Contents file.
\newwrite\tocfile 

% New Index file
\newwrite\indexfile 

% New line character.
\newlinechar = `\^^J

% Open the index file.
\immediate\openout\indexfile = \jobname index.idx

\def\index#1{%
   #1%
   \immediate\write\indexfile
   {#1,\folio}%
   }%

% Begin table of contents processing. Write contents to file \tocfile.
\def\maketoc{%

   % Open the table of contents file.
   \immediate\openin\tocfile = \jobname toc.tex

   \immediate\ifeof\tocfile {\message {No table of contents file 
   \jobname toc.tex . Run TeX again to produce the table of contents.}} 
   \else \closein \tocfile \input \jobname toc.tex \fi

   % Open the table of contents file.
   \immediate\openout\tocfile = \jobname toc.tex

   \immediate \write \tocfile 
   {
      {\noexpand\bxmplbxt {\noexpand\centerline {Contents}}}
      ^^J
      \noexpand\vskip .2in
      ^^J
      \noexpand\hrule height 1pt
      ^^J
      \noexpand\vskip .2in
      ^^J
      \noexpand\begingroup
      ^^J
      ^^J
      ^^J
      }
   }

% End table of contents processing. Write contents to file \tocfile.
\def\endbook{%
   \immediate \write \tocfile 
   {\noexpand\endgroup
      ^^J
      \noexpand\vfil
      ^^J
      \noexpand\supereject
      ^^J
   }%
   \closeout \tocfile
   \closeout \indexfile
   }%

% Set the footline for the document
% The title page sets pageno to 0 so that no page number appears.
% The preface sets pageno to -1 to get roman numeral page numbers.
% Chapter I sets page number to 1.
\footline={\ifnum \pageno = 0 \hfil\else \tenit \hfil\folio\hfil\fi}
\headline={\ifnum \count6 = 0 \hfil\else \hfil {\header \firstmark}\fi}

% A section header. 
\def\section#1{\vskip .25 in 
  \goodbreak 
  \leftline{\xmplbx #1}\mark{#1}
  \count6=\count3 % Enable running header once a section has appeared.
  \immediate\write \tocfile 
  {\vskip .025in 
  \noexpand\line{{\noexpand \hskip .2in {\sectioncont #1}\/}
  \noexpand\dotting \number\pageno}
  \vskip .025in}
  \smallskip
  }%

% A subsection header. 
\def\subsection#1{\vskip .15 in 
  \goodbreak 
  \leftline{\sxmplbx #1}\mark{#1} 
  \smallskip
  }%

% Chapter macro. If at first chapter set \pageno to 1.
\count3=0 % Chapter Number.
\count6=0 % Chapter Number (after first chapter page; otherwise 0)
\def\chapter#1{%
   \count1=0   % Figures
   \count4=0   % Exercises
   \count5=0   % Examples
   \vfil       % Fill out last page.
   \supereject % End the page.
   \mark{}
   \count6=0   % Chapter number (after first page)
   \global\advance\count3 by 1
   {\bxmplbx {\centerline {Chapter \number\count3: #1}}}
   \vskip .2in
   \hrule height 1pt
   \vskip .2in
   \ifnum \count3 = 1 \pageno=1 \fi
   \immediate\write \tocfile 
   {\vskip .1in
   \noexpand\line{{\noexpand{\chap Chapter \number\count3:} {\chapcont #1}\/}
   \noexpand\dotting \number\pageno}
   \vskip .1in}
   }%

\newdimen\exerindent \exerindent=6pc
\def\exercise#1{%
   \global\advance\count4 by 1
   \goodbreak
   {\noindent \leftskip = \exerindent
      \ifnum \number\count3 > 0
      \llap{\hbox to \exerindent{\bf Exercise \number\count3.\number\count4:\hfil}}%
      \else
      \llap{\hbox to \exerindent{\bf Exercise \number\count4:\hfil}}%
      \fi
      #1\par
      }%
   }%

\newdimen\examindent \examindent=6pc
\def\example#1{%
\global\advance\count5 by 1
\goodbreak
{\noindent \leftskip = \examindent
\ifnum \number\count3 > 0
\llap{\hbox to \examindent{\bf Example \number\count3.\number\count5:\hfil}}%
\else
\llap{\hbox to \examindent{\bf Example \number\count5:\hfil}}%
\fi
#1\par
}%
}%

\def\createindex{%
\vfil     % Fill the last page.
\supereject % End last page.
{\centerline {\bxmplbx Index}}
\vskip .2in
\hrule height 1pt
\vskip .2in
}

% Appendix macro.
\def\appendix#1#2{%
\count1=0 % Figure Number
\count3=0 % Chapter Number
\count4=0 % Exercises
\count5=0 % Examples
\vfil     % Fill the last page.
\supereject % End last page.
{\leftline {\bxmplbx  {#1:\ #2}}}
\vskip .2in
\hrule height 1pt
\vskip .2in
\immediate\write \tocfile 
{\vskip .1in
\noexpand\line{{\noexpand {\chap #1:\/}}\ {\chapcont #2}%
\noexpand\dotting \number\pageno}
\vskip .1in}
}%

% A title macro. First parameter is the title of the article.
% My name and the current time appear in the title section.
\def\titledate#1#2{%
{\bxmplbx {\centerline {#1}}}
\vskip .2in
  {\leftskip = 3in%
  \par\noindent\llap{\hbox to 3in{\hfil {\bf Author:} }}{\it R. Scott
  McIntire}\par\noindent {\it Financial Engineering}\par\noindent 
  {\it ITG, Inc.}\par%
  }%
\vskip .1in
 {\leftskip = 3in%
  \par\noindent\llap{\hbox to 3in{\hfil {\bf Date:} }}#2 \par%
  }%
\vskip .25in
\hrule height 1pt
\vskip .1in
}

% A title macro. First parameter is the title of the article.
% Second parameter is who; third parameter is date.
\def\mantitle#1#2#3{%
{\bxmplbx {\centerline {#1}}}
\vskip .2in
  {\leftskip = 2.5in%
  \par\noindent\llap{\hbox to 2.5in{\hfil {\bf Author:} }}{\it #2}
  \par%
  }%
\vskip .1in
 {\leftskip = 2.5in%
  \par\noindent\llap{\hbox to 2.5in{\hfil {\bf Date:} }}#3\par%
  }%
\vskip .25in
\hrule height 1pt
\vskip .1in
}

% Produce a fraction #1/#2.
\def\frac#1#2{{{\scriptstyle #1} \over {\scriptstyle #2}}}

% Produce a fraction #1/#2 in display style.
\def\fracd#1#2{{{\displaystyle #1} \over {\displaystyle #2}}}

% Create an overbrace over expression #1 with text #2 above.
\def\obrace#1#2{%
\mathop{\overbrace{#1}}\limits^{\hbox{\small #2}}
}%

% Figure macro for eps files. #1 - eps file, #2 - caption text.
\count1=0 % Figure Number
\def\epsfigure#1#2{%
\global\advance\count1 by 1
  \bigskip
  \goodbreak
  \vbox{\hfil \epsfbox{#1} \hfil
      \vskip 0pt
      \ifnum \number\count3 > 0
{ \noindent \leftskip=\leftfigureskip \rightskip=\rightfigureskip 
       \llap{Figure \number\count3.\number\count1:\quad}#2\par}
      \else
{ \noindent \leftskip=\leftfigureskip \rightskip=\rightfigureskip
       \llap{Figure \number\count1:\quad}#2\par}
      \fi
  \bigskip}}

% Item macro.
\def\item#1{%
{\leftskip = 4pc \parskip=0pt%
\par\noindent\llap{\hbox to 4pc{\hfil $\bf \bullet$ \hfil}}#1\par%
}%
}

% Begin enumeration marker.
\def\beginEnum{%
\vskip 0.3\parskip
\count2=0
}

% Enumerate text.
\def\enum#1{%
  \global\advance\count2 by 1
  {\leftskip=2.5pc \parskip=0pt%
  \par\noindent\llap{\hbox to 2pc{\hfil $\bf \number\count2$. \hfil}}#1\par%
  }%
}

% End enumeration marker.
\def\endEnum{%
   \count2=0
   \vskip -.75\parskip
}

% Dotting used in the table of contents.
\def\dotting{\leaders\hbox to 1em{\hfil.\hfil}\hfil}

